\chapter{2006年湖南卷（文）}
\section{单选题}
\begin{problem}
函数 \(y=\sqrt{\log_{2}x}\) 的定义域是（\quad）

\choices
    {\((0,1]\)}
    {\((0,+\infty)\)}
    {\((1,+\infty)\)}
    {\([1,+\infty)\)}
\end{problem}

\begin{answer}
D
\end{answer}

\begin{solution}
根式有意义要求 \(\log_{2}x\geqslant0\)，同时对数的真数必须满足 \(x>0\). 因为底数 \(2>1\)，所以 \(\log_{2}x\geqslant0\) 等价于 \(x\geqslant2^0=1\). 故函数的定义域为 \([1,+\infty)\)，选 D.
\end{solution}

\begin{problem}
已知向量 \(\bs{a}=(2,t)\)，\(\bs{b}=(1,2)\)，若 \(t=t_1\) 时，\(\bs{a}\parallel\bs{b}\)；\(t=t_2\) 时，\(\bs{a}\perp\bs{b}\)，则（\quad）

\choices
    {\(t_1=-4,t_2=-1\)}
    {\(t_1=-4\)，\(t_2=1\)}
    {\(t_1=4\)，\(t_2=-1\)}
    {\(t_1=4\)，\(t_2=1\)}
\end{problem}

\begin{answer}
C
\end{answer}

\begin{solution}
向量平行时对应坐标成比例，故 \(2\times2-t\times1=0\)，从而 \(t_1=4\). 向量垂直时数量积为零，故 \(\bs{a}\cdot\bs{b}=2\times1+2t=0\)，从而 \(t_2=-1\). 因此选 C.
\end{solution}

\begin{problem}
若 \((ax-1)^5\) 的展开式中 \(x^3\) 的系数是 \(80\)，则实数 \(a\) 的值是（\quad）

\choices
    {\(-2\)}
    {\(2\sqrt{2}\)}
    {\(\sqrt[3]{4}\)}
    {\(2\)}
\end{problem}

\begin{answer}
D
\end{answer}

\begin{solution}
\(x^3\) 项必须从五个因式中取出三个 \(ax\) 和两个 \(-1\)，所以其系数为 \(\mathrm C_5^3a^3(-1)^2=10a^3\). 由题意得 \(10a^3=80\)，即 \(a^3=8\)，因此 \(a=2\)，选 D.
\end{solution}

\begin{problem}
过半径为 \(2\) 的球 \(O\) 表面上一点 \(A\) 作球 \(O\) 的截面，若 \(OA\) 与该截面所成的角是 \(60^{\circ}\)，则该截面的面积是（\quad）

\choices
    {\(\pi\)}
    {\(2\pi\)}
    {\(3\pi\)}
    {\(2\sqrt{3}\pi\)}
\end{problem}

\begin{answer}
A
\end{answer}

\begin{solution}
设截面圆心为 \(H\). 球心到截面的垂线 \(OH\) 垂直于截面，而 \(AH\) 是半径 \(2\) 的球在截面内的投影，故 \(\triangle OAH\) 为直角三角形，且 \(\angle OAH=60^{\circ}\). 于是截面圆半径 \(AH=OA\cos60^{\circ}=2\times\frac12=1\)，截面面积为 \(\pi AH^2=\pi\)，选 A.
\end{solution}

\begin{problem}
“\(a=1\)”是“函数 \(f(x)=|x-a|\) 在区间 \([1,+\infty)\) 上为增函数”的（\quad）

\choices
    {充分不必要条件}
    {必要不充分条件}
    {充要条件}
    {既不充分也不必要条件}
\end{problem}

\begin{answer}
A
\end{answer}

\begin{solution}
若 \(a=1\)，则对 \(x\geqslant1\)，有 \(f(x)=|x-1|=x-1\)，所以函数在 \([1,+\infty)\) 上为增函数. 反过来，若 \(a\leqslant1\)，则同样有 \(f(x)=x-a\)，函数仍在该区间上为增函数；例如 \(a=0\) 时条件成立但 \(a\neq1\). 因此“\(a=1\)”是充分不必要条件，选 A.
\end{solution}

\begin{problem}
在数字 \(1\)，\(2\)，\(3\) 与符号 \(+\)，\(-\) 五个元素的所有全排列中，任意两个数字都不相邻的全排列个数是（\quad）

\choices
    {\(6\)}
    {\(12\)}
    {\(18\)}
    {\(24\)}
\end{problem}

\begin{answer}
B
\end{answer}

\begin{solution}
先排列三个数字，有 \(3!\) 种. 为使任意两个数字不相邻，两个符号必须分别放在三个数字之间的两个空隙中；两个符号还可以交换位置，有 \(2!\) 种. 因此全排列个数为 \(3!\times2!=12\)，选 B.
\end{solution}

\begin{problem}
圆 \(x^2+y^2-4x-4y-10=0\) 上的点到直线 \(x+y-14=0\) 的最大距离与最小距离的差是（\quad）

\choices
    {\(36\)}
    {\(18\)}
    {\(6\sqrt{2}\)}
    {\(5\sqrt{2}\)}
\end{problem}

\begin{answer}
C
\end{answer}

\begin{solution}
配方得 \((x-2)^2+(y-2)^2=18\)，所以圆心为 \(C(2,2)\)，半径为 \(r=3\sqrt{2}\). 圆心到直线的距离为 \(d=\dfrac{|2+2-14|}{\sqrt{1^2+1^2}}=5\sqrt{2}\)，且 \(d>r\). 因此圆上点到该直线的最大、最小距离分别为 \(d+r\) 和 \(d-r\)，其差为 \(2r=6\sqrt{2}\)，选 C.
\end{solution}

\begin{problem}
设点 \(P\) 是函数 \(f(x)=\sin\omega x\) 的图象 \(C\) 的一个对称中心，若点 \(P\) 到图象 \(C\) 的对称轴上的距离的最小值是 \(\frac{\pi}{4}\)，则 \(f(x)\) 的最小正周期是（\quad）

\choices
    {\(2\pi\)}
    {\(\pi\)}
    {\(\frac{\pi}{2}\)}
    {\(\frac{\pi}{4}\)}
\end{problem}

\begin{answer}
B
\end{answer}

\begin{solution}
函数 \(y=\sin\omega x\) 的对称中心为 \(\left(\dfrac{k\pi}{\omega},0\right)\)，对称轴为 \(x=\dfrac{(2j+1)\pi}{2\omega}\)，其中 \(k\)，\(j\in\Z\). 相邻对称中心与对称轴的横坐标距离为 \(\dfrac{\pi}{2|\omega|}\)，题设给出 \(\dfrac{\pi}{2|\omega|}=\dfrac{\pi}{4}\)，所以 \(|\omega|=2\). 故最小正周期为 \(T=\dfrac{2\pi}{|\omega|}=\pi\)，选 B.
\end{solution}

\begin{problem}
过双曲线 \(M:x^2-\dfrac{y^2}{b^2}=1\) 的左顶点 \(A\) 作斜率为 \(1\) 的直线 \(l\)，若 \(l\) 与双曲线 \(M\) 的两条渐近线分别相交于点 \(B\)，\(C\)，且 \(|AB|=|BC|\)，则双曲线 \(M\) 的离心率是（\quad）

\choices
    {\(\frac{\sqrt{5}}{2}\)}
    {\(\frac{\sqrt{10}}{3}\)}
    {\(\sqrt{5}\)}
    {\(\sqrt{10}\)}
\end{problem}

\begin{answer}
D
\end{answer}

\begin{solution}
双曲线的左顶点为 \(A(-1,0)\)，渐近线为 \(y=\pm bx\)，直线 \(l\) 的方程为 \(y=x+1\). 令直线上的点表示为 \((-1+t,t)\)，则与 \(y=-bx\) 和 \(y=bx\) 的交点参数分别为 \(t_1=\dfrac{b}{b+1}\) 和 \(t_2=\dfrac{b}{b-1}\). 当 \(0<b<1\) 时两个参数异号，较远的两点距离大于从 \(A\) 到任一交点的距离，不能满足 \(|AB|=|BC|\)；当 \(b=1\) 时直线与一条渐近线平行. 因此必须有 \(b>1\). 此时 \(0<t_1<1<t_2\)，满足距离条件的标号顺序只能是 \(A\)、参数为 \(t_1\) 的点 \(B\)、参数为 \(t_2\) 的点 \(C\). 由于直线参数每增加 \(1\)，点沿直线移动 \(\sqrt{2}\)，故
\(|AB|=t_1\sqrt{2}\)，\(|BC|=(t_2-t_1)\sqrt{2}\).
由 \(|AB|=|BC|\) 得 \(t_2=2t_1\)，即
\[
\frac{b}{b-1}=\frac{2b}{b+1}
\]
约去 \(b>0\) 得 \(b+1=2b-2\)，所以 \(b=3\). 双曲线的实半轴为 \(1\)，故离心率 \(e=\sqrt{1+b^2}=\sqrt{10}\)，选 D.
\end{solution}

\begin{problem}
如图，\(OM\parallel AB\)，点 \(P\) 由射线 \(OM\)、线段 \(OB\) 及 \(AB\) 的延长线围成的阴影区域内（不含边界），且 \(\overrightarrow{OP}=x\overrightarrow{OA}+y\overrightarrow{OB}\)，则实数对 \((x,y)\) 可以是（\quad）

\bitmapfigure[max width=0.40\linewidth]{img/2006/hunan_liberal/q10_fig1.png}

\choices
    {\(\left(\frac{1}{4},\frac{3}{4}\right)\)}
    {\(\left(-\frac{2}{3},\frac{2}{3}\right)\)}
    {\(\left(-\frac{1}{4},\frac{3}{4}\right)\)}
    {\(\left(-\frac{1}{5},\frac{7}{5}\right)\)}
\end{problem}

\begin{answer}
C
\end{answer}

\begin{solution}
设 \(\overrightarrow{OA}=\bs{a}\)，\(\overrightarrow{OB}=\bs{b}\). 阴影区域内的点可以唯一表示为
\(\overrightarrow{OP}=\lambda_1\overrightarrow{OB}+\lambda_2\overrightarrow{AB}\)，\(0<\lambda_1<1\)，\(\lambda_2>0\).
又因为 \(\overrightarrow{AB}=\overrightarrow{OB}-\overrightarrow{OA}\)，所以
\[
\overrightarrow{OP}=-\lambda_2\overrightarrow{OA}+\left(\lambda_1+\lambda_2\right)\overrightarrow{OB}
\]
比较 \(\overrightarrow{OP}=x\overrightarrow{OA}+y\overrightarrow{OB}\) 的系数，得 \(x=-\lambda_2<0\)，且 \(x+y=\lambda_1\in(0,1)\). A 的第一坐标为正，B 的 \(x+y=0\) 在边界上，D 的 \(x+y=\frac65>1\)；只有 C 满足 \(x=-\frac14<0\) 且 \(x+y=\frac12\in(0,1)\). 故选 C.
\end{solution}

\section{填空题}
\begin{problem}
若数列 \(\{a_n\}\) 满足：\(a_1=1\)，\(a_{n+1}=2a_n\)，\(n=1\)，\(2\)，\(3\)，\(\cdots\)，则 \(a_1+a_2+\cdots+a_n=\)\fillinblank{}.
\end{problem}

\begin{answer}
\(2^n-1\)
\end{answer}

\begin{solution}
由 \(a_{n+1}=2a_n\) 和 \(a_1=1\) 得 \(a_k=2^{k-1}\). 因此
\[
a_1+a_2+\cdots+a_n=\sum_{k=1}^{n}2^{k-1}=\frac{2^n-1}{2-1}=2^n-1
\]
\end{solution}

\begin{problem}
某高校有甲，乙两个数学建模兴趣班. 其中甲班 \(40\text{ 人}\)，乙班 \(50\text{ 人}\). 现分析两个班的一次考试成绩，算得甲班的平均成绩是 \(90\text{ 分}\)，乙班的平均成绩是 \(81\text{ 分}\)，则该校数学建模兴趣班的平均成绩是\fillinblank{}\(\text{ 分}\).
\end{problem}

\begin{answer}
\(85\)
\end{answer}

\begin{solution}
甲、乙两班的总成绩分别为 \(40\times90\) 和 \(50\times81\)，总人数为 \(40+50=90\). 所以全班平均成绩为
\[
\frac{40\times90+50\times81}{40+50}=\frac{7650}{90}=85
\]
\end{solution}

\begin{problem}
已知 \(\begin{cases}
x\geqslant1,\\
x-y+1\leqslant0,\\
2x-y-2\leqslant0
\end{cases}\)，则 \(x^2+y^2\) 的最小值是\fillinblank{}.
\end{problem}

\begin{answer}
\(5\)
\end{answer}

\begin{solution}
约束条件等价于 \(x\geqslant1\)、\(y\geqslant x+1\)、\(y\geqslant2x-2\). 当 \(1\leqslant x\leqslant3\) 时，\(x+1\geqslant2x-2\)，且 \(y\geqslant x+1>0\)，所以
\[
x^2+y^2\geqslant x^2+(x+1)^2=2x^2+2x+1
\]
该式在 \([1,3]\) 上递增，最小值在 \(x=1\)，\(y=2\) 处取得，为 \(5\). 当 \(x\geqslant3\) 时，\(2x-2\geqslant x+1\)，所以
\[
x^2+y^2\geqslant x^2+(2x-2)^2=5x^2-8x+4\geqslant25
\]
比较两种情形，原式的最小值为 \(5\).
\end{solution}

\begin{problem}
过三棱柱 \(ABC-A_1B_1C_1\) 的任意两条棱的中点作直线，其中与平面 \(ABB_1A_1\) 平行的直线共有\fillinblank{}\(\text{条}\).
\end{problem}

\begin{answer}
\(6\)
\end{answer}

\begin{solution}
取 \(A\) 为原点，令 \(\bs{u}=\overrightarrow{AB}\)，\(\bs{v}=\overrightarrow{AC}\)，\(\bs{w}=\overrightarrow{AA_1}\). 平面 \(ABB_1A_1\) 内的点不含 \(\bs{v}\) 方向的分量. 九条棱的中点按其 \(\bs{v}\) 分量分为三组：\(AB\)，\(A_1B_1\)，\(AA_1\)，\(BB_1\) 的中点分量为 \(0\)；\(AC\)，\(BC\)，\(A_1C_1\)，\(B_1C_1\) 的中点分量为 \(\frac12\)；\(CC_1\) 的中点分量为 \(1\).

第二组的四个中点共面于一个与平面 \(ABB_1A_1\) 平行且不相交的平面，任取其中两个中点所作直线都与平面 \(ABB_1A_1\) 平行，因此共有 \(\mathrm C_4^2=6\) 条. 第一组中点连线落在原平面内，不属于与平面不相交的平行直线；不同组之间的连线含有 \(\bs{v}\) 方向分量，也不平行. 故共有 \(6\) 条.
\end{solution}

\begin{problem}
若 \(f(x)=a\sin\left(x+\frac{\pi}{4}\right)+3\sin\left(x-\frac{\pi}{4}\right)\) 是偶函数，则 \(a=\)\fillinblank{}.
\end{problem}

\begin{answer}
\(-3\)
\end{answer}

\begin{solution}
由和差公式，
\[
f(x)=\frac{a+3}{\sqrt{2}}\sin x+\frac{a-3}{\sqrt{2}}\cos x
\]
其中 \(\sin x\) 是奇函数，\(\cos x\) 是偶函数. 要使 \(f(x)\) 为偶函数，奇函数项的系数必须为零，即 \(\dfrac{a+3}{\sqrt{2}}=0\)，所以 \(a=-3\).
\end{solution}

\section{解答题}
\begin{problem}
已知 \(\sqrt{3}\sin\theta-\dfrac{\sin\left(\frac{\pi}{2}-2\theta\right)}{\cos\left(\pi+\theta\right)}\cdot\cos\theta=1\)，\(\theta\in(0,\pi)\)，求 \(\theta\) 的值.
\end{problem}

\begin{solution}
原式有意义要求 \(\cos(\pi+\theta)=-\cos\theta\neq0\). 由诱导公式，\(\sin\left(\frac{\pi}{2}-2\theta\right)=\cos2\theta\)，所以原方程化为
\[
\sqrt{3}\sin\theta+\cos2\theta=1
\]
再用 \(\cos2\theta=1-2\sin^2\theta\)，得 \(\sin\theta\left(\sqrt{3}-2\sin\theta\right)=0\). 由于 \(0<\theta<\pi\)，有 \(\sin\theta>0\)，所以 \(\sin\theta=\frac{\sqrt{3}}{2}\). 于是 \(\theta=\frac{\pi}{3}\) 或 \(\theta=\frac{2\pi}{3}\)，且这两个值都满足 \(\cos\theta\neq0\).
\end{solution}

\begin{problem}
某安全生产监督部门对 \(5\text{ 家}\)小型煤矿进行安全检查（简称安检）. 若安检不合格，则必须整改. 若整改后经复查仍不合格，则强制关闭. 设每家煤矿安检是否合格是相互独立的，且每家煤矿整改前安检合格的概率是 \(0.5\)，整改后安检合格的概率是 \(0.8\)，计算（结果精确到 \(0.01\)）：

\begin{enumerate}
\item 恰好有两家煤矿必须整改的概率；
\item 某煤矿不被关闭的概率；
\item 至少关闭一家煤矿的概率.
\end{enumerate}
\end{problem}

\begin{solution}
\begin{enumerate}
\item 每家煤矿整改前安检不合格的概率为 \(1-0.5=0.5\)，所以恰好有两家必须整改的概率为
\[
\mathrm C_5^2(0.5)^2(0.5)^3=\frac{5}{16}=0.3125\approx0.31
\]
\item 某一家煤矿不被关闭包括两种互斥情形：整改前安检合格，或整改前不合格但整改后复查合格. 因此概率为 \(0.5+0.5\times0.8=0.9\approx0.90\).
\item 一家煤矿被关闭的概率为 \(0.5\times(1-0.8)=0.1\)，所以一家煤矿不被关闭的概率为 \(0.9\). 由各家煤矿相互独立，五家都不被关闭的概率为 \((0.9)^5\)，故至少关闭一家煤矿的概率为
\[
1-(0.9)^5=0.40951\cdots\approx0.41
\]
\end{enumerate}
\end{solution}

\begin{problem}
如图，已知两个正四棱锥 \(P-ABCD\) 与 \(Q-ABCD\) 的高都是 \(2\)，\(AB=4\).

\begin{enumerate}
\item 证明：\(PQ\perp\) 平面 \(ABCD\)；
\item 求异面直线 \(AQ\) 与 \(PB\) 所成的角；
\item 求点 \(P\) 到平面 \(QAD\) 的距离.
\end{enumerate}

\bitmapfigure[max width=0.42\linewidth]{img/2006/hunan_liberal/q18_fig1.png}
\end{problem}

\begin{solution}
\begin{enumerate}
\item 设正方形 \(ABCD\) 的中心为 \(O\). 正四棱锥的高垂直于底面且交于底面中心，所以 \(PO\perp\) 平面 \(ABCD\)，\(QO\perp\) 平面 \(ABCD\). 过平面内同一点 \(O\) 的两条垂线重合，故 \(P\)，\(O\)，\(Q\) 共线，从而 \(PQ\perp\) 平面 \(ABCD\).
\item 建立直角坐标系，取
\(A(-2,-2,0)\)，\(B(2,-2,0)\)，\(D(-2,2,0)\)，\(P(0,0,2)\)，\(Q(0,0,-2)\).
则 \(\overrightarrow{AQ}=(2,2,-2)\)，\(\overrightarrow{PB}=(2,-2,-2)\)，所以
\(\overrightarrow{AQ}\cdot\overrightarrow{PB}=4\)，\(|\overrightarrow{AQ}|=|\overrightarrow{PB}|=2\sqrt{3}\).
设异面直线所成的锐角为 \(\alpha\)，则 \(\cos\alpha=\dfrac{|4|}{(2\sqrt{3})(2\sqrt{3})}=\frac13\)，故所求角为 \(\arccos\dfrac13\).
\item 由 \(\overrightarrow{AD}=(0,4,0)\) 和 \(\overrightarrow{AQ}=(2,2,-2)\) 可取平面 \(QAD\) 的一个法向量为 \((1,0,1)\)，故其方程为 \(x+z+2=0\). 点 \(P(0,0,2)\) 到该平面的距离为
\[
\frac{|0+2+2|}{\sqrt{1^2+1^2}}=2\sqrt{2}
\]
\end{enumerate}
\end{solution}

\begin{problem}
已知函数 \(f(x)=ax^3-3x^2+1-\dfrac{3}{a}\).

\begin{enumerate}
\item 讨论函数 \(f(x)\) 的单调性；
\item 若曲线 \(y=f(x)\) 上两点 \(A\)，\(B\) 处的切线都与 \(y\) 轴垂直，且线段 \(AB\) 与 \(x\) 轴有公共点，求实数 \(a\) 的取值范围.
\end{enumerate}
\end{problem}

\begin{solution}
\begin{enumerate}
\item 因为函数有定义，所以 \(a\neq0\)，且
\[
f'(x)=3ax^2-6x=3x(ax-2)
\]
当 \(a>0\) 时，\(0<\dfrac{2}{a}\). 由导数符号可知，函数在 \((-\infty,0)\) 和 \(\left(\dfrac{2}{a},+\infty\right)\) 上递增，在 \(\left(0,\dfrac{2}{a}\right)\) 上递减. 当 \(a<0\) 时，\(\dfrac{2}{a}<0\)，函数在 \(\left(-\infty,\dfrac{2}{a}\right)\) 和 \((0,+\infty)\) 上递减，在 \(\left(\dfrac{2}{a},0\right)\) 上递增.
\item 切线与 \(y\) 轴垂直等价于切线水平，即 \(f'(x)=0\). 两点的横坐标为 \(0\) 和 \(\dfrac{2}{a}\)，对应纵坐标为
\(f(0)=1-\frac{3}{a}\)，\(f\left(\frac{2}{a}\right)=1-\frac{3}{a}-\frac{4}{a^2}=\frac{(a-4)(a+1)}{a^2}\).
线段 \(AB\) 与 \(x\) 轴有公共点，当且仅当两个端点的纵坐标异号或至少一个为零，即
\[
\left(1-\frac{3}{a}\right)\frac{(a-4)(a+1)}{a^2}\leqslant0
\]
也就是
\(\frac{(a-3)(a-4)(a+1)}{a^3}\leqslant0\)，\(a\neq0\).
在区间 \((-\infty,-1)\)、\((-1,0)\)、\((0,3)\)、\((3,4)\)、\((4,+\infty)\) 内，上式的符号依次为正、负、正、负、正；在 \(a=-1\)，\(3\)，\(4\) 时上式等于 \(0\)，而 \(a=0\) 不在定义域内. 因此 \(a\in[-1,0)\cup[3,4]\).
\end{enumerate}
\end{solution}

\begin{problem}
在 \(m\)（\(m\geqslant2\)）个不同数的排列 \(P_1P_2\cdots P_m\) 中，若 \(1\leqslant i<j\leqslant m\) 时 \(P_i>P_j\)（即前面某数大于后面某数），则称 \(P_i\) 与 \(P_j\) 构成一个逆序. 一个排列的全部逆序的总数称为该排列的逆序数. 记排列 \((n+1)n(n-1)\cdots321\) 的逆序数为 \(a_n\)，如排列 \(21\) 的逆序数 \(a_1=1\)，排列 \(321\) 的逆序数 \(a_2=3\)，排列 \(4321\) 的逆序数 \(a_3=6\).

\begin{enumerate}
\item 求 \(a_4\)，\(a_5\)，并写出 \(a_n\) 的表达式；
\item 令 \(b_n=\dfrac{a_n}{a_{n+1}}+\dfrac{a_{n+1}}{a_n}\)，证明：\(2n<b_1+b_2+\cdots+b_n<2n+3\)，\(n=1\)，\(2\)，\(\cdots\).
\end{enumerate}
\end{problem}

\begin{solution}
\begin{enumerate}
\item 完全逆序排列中，任取其中两个数，较大的数都在前面，因而这两个数恰好构成一个逆序. 共有 \(n+1\) 个数，所以
\[
a_n=\mathrm C_{n+1}^2=\frac{n(n+1)}{2}
\]
从而 \(a_4=\mathrm C_5^2=10\)，\(a_5=\mathrm C_6^2=15\).
\item 由上式得
\[
b_n=\frac{n}{n+2}+\frac{n+2}{n}=2+\frac{4}{n(n+2)}=2+2\left(\frac1n-\frac1{n+2}\right)
\]
因此
\[
\begin{gathered}
b_1+b_2+\cdots+b_n
 =2n+2\sum_{k=1}^{n}\left(\frac1k-\frac1{k+2}\right)\\
 =2n+2\left(1+\frac12-\frac1{n+1}-\frac1{n+2}\right)\\
 =2n+3-\frac2{n+1}-\frac2{n+2}.
\end{gathered}
\]
由于 \(\dfrac2{n+1}+\dfrac2{n+2}>0\)，上式小于 \(2n+3\)；又因 \(n\geqslant1\) 时 \(\dfrac2{n+1}+\dfrac2{n+2}<3\)，上式大于 \(2n\). 故不等式成立.
\end{enumerate}
\end{solution}

\begin{problem}
已知椭圆 \(C_1:\dfrac{x^2}{4}+\dfrac{y^2}{3}=1\)，抛物线 \(C_2:(y-m)^2=2px\)（\(p>0\)），且 \(C_1\)，\(C_2\) 的公共弦 \(AB\) 过椭圆 \(C_1\) 的右焦点.

\begin{enumerate}
\item 当 \(AB\perp x\) 轴时，求 \(m\)，\(p\) 的值，并判断抛物线 \(C_2\) 的焦点是否在直线 \(AB\) 上；
\item 若 \(p=\frac43\) 且抛物线 \(C_2\) 的焦点在直线 \(AB\) 上，求 \(m\) 的值及直线 \(AB\) 的方程.
\end{enumerate}
\end{problem}

\begin{solution}
\begin{enumerate}
\item 椭圆的半长轴、半短轴分别为 \(2\)、\(\sqrt3\)，所以半焦距为 \(c=\sqrt{2^2-(\sqrt3)^2}=1\)，右焦点为 \(F(1,0)\). 因为 \(AB\perp x\) 轴且过 \(F\)，所以 \(AB\) 的方程为 \(x=1\). 代入椭圆方程得公共弦端点的纵坐标为 \(y=\pm\frac32\). 两端点在抛物线上，故
\(\left(\frac32-m\right)^2=2p\)，\(\left(-\frac32-m\right)^2=2p\).
两式相减得 \(m=0\)，代回得 \(2p=\frac94\)，所以 \(p=\frac98\). 此时抛物线焦点为 \(\left(\frac{p}{2},m\right)=\left(\frac9{16},0\right)\)，横坐标不等于 \(1\)，故焦点不在直线 \(AB\) 上.
\item 当 \(p=\frac43\) 时，抛物线焦点为 \(\left(\frac23,m\right)\). 公共弦不可能是竖直线：若为 \(x=1\)，由第一问应有 \(p=\frac98\)；也不可能是水平线，因为水平线 \(y=0\) 与椭圆的两个交点横坐标为 \(\pm2\)，不可能同时满足抛物线方程. 因此设 \(AB\) 为 \(y=k(x-1)\)，其中 \(k\neq0\). 焦点在该直线上给出 \(m=-\frac{k}{3}\).

将直线代入椭圆方程，得交点横坐标满足
\[
(3+4k^2)x^2-8k^2x+4k^2-12=0
\]
将 \(m=-\frac{k}{3}\) 和 \(p=\frac43\) 代入抛物线方程，得同一对交点的横坐标还满足
\[
9k^2x^2-(12k^2+24)x+4k^2=0
\]
两式有相同的两个不同实根，所以对应系数成比例. 比较二次项和常数项，得
\[
9k^2(4k^2-12)=4k^2(3+4k^2)
\]
因 \(k\neq0\)，化简得 \(20k^2=120\)，即 \(k^2=6\). 于是
\(k=\sqrt6\Rightarrow m=-\frac{\sqrt6}{3}\)，\(AB:y=\sqrt6(x-1);\)
\(k=-\sqrt6\Rightarrow m=\frac{\sqrt6}{3}\)，\(AB:y=-\sqrt6(x-1)\).
\end{enumerate}
\end{solution}
